% !PRINTED 484
\weberpage{484}
%<<CONT
dicts the assumption that the values (3) are different from one another. Permutation
groups $P$ which have this property, namely that the permuted
elements can be decomposed into rows $A, B \dots S$ of equally
many elements in such a way that no permutation of $P$ carries
two elements of the same row into elements of different rows,
are called \gesperrt{imprimitive}. The individual rows
$A, B \dots S$ are called the \gesperrt{systems of imprimitivity}.
Permutation groups for which such a decomposition of the elements
is not possible are called \gesperrt{primitive groups}.

It may very well happen that a group is imprimitive in more
than one way, that it possesses entirely different systems of
imprimitivity. Thus, for example, the group of cyclic permutations
of six elements $(1, 2, 3, 4, 5, 6)$ is imprimitive in a twofold
manner and has the systems of imprimitivity

\[
A = 1, 3, 5, \quad B = 2, 4, 6
\]
and
\[
A = 1, 4, \quad B = 2, 5, \quad C = 3, 6.
\]

Now that the concept of primitive and imprimitive groups has
thus been established, we can give the following expression to
what was proved above:

1. \gesperrt{An imprimitive field has an imprimitive
group.}

From the imprimitivity of the group there results an important
consequence for the imprimitive fields.

We will denote by
\[
\psi(x, x_1, \dots x_{r-1})
\]
a rational \gesperrt{symmetric function} of the $r$ variables
$x, x_1 \dots x_{r-1}$, and set

\weq{4}{
\begin{aligned}
\omega &= \psi(\alpha, \alpha_1 \dots \alpha_{r-1}) \\
\omega_1 &= \psi(\beta, \beta_1 \dots \beta_{r-1}) \\
& \cdot \ \cdot \ \cdot \ \cdot \ \cdot \ \cdot \ \cdot \ \cdot \ \cdot \ \cdot \\
\omega_{s-1} &= \psi(\sigma, \sigma_1 \dots \sigma_{r-1});
\end{aligned}
}
then it is to be proved that $\omega$ can be expressed rationally
in terms of $\Theta$, that is, is contained in the field $\Omega(\Theta)$. For
we have
\weq{5}{
\varphi(t)\left(\dfrac{\omega}{t - \Theta} + \dfrac{\omega_1}{t - \Theta_1} + \cdots + \dfrac{\omega_{s-1}}{t - \Theta_{s-1}}\right) = \Phi(t)
}
%CONT>>
